\section{Periodic Point Detection (Feature Finder)}
\label{sec:feature-finder}

Given a viewer location and a search radius, \FractalShark{} can automatically
locate a candidate center of a nearby mini-Mandelbrot copy by searching for a
parameter $c$ whose critical orbit returns to zero.  The search radius limits
the region of interest; Newton iteration is a local solver and does not
guarantee selection of the nearest center.  This capability is
implemented by the \code{FeatureFinder} class (files
\code{FeatureFinder.h}/\code{FeatureFinder.cpp}), whose design is heavily
influenced by the periodic-point finder in the Imagina fractal viewer \cite{Imagina}.  The
algorithm proceeds in two phases:
\begin{enumerate}
\item \textbf{Phase~A (T-space).}
  Detect the candidate period $p$ and perform a coarse Newton iteration using
  the renderer's native floating-point type~$T$ (which may be \code{float},
  \code{double}, or \code{HDRFloat}).
\item \textbf{Phase~B (MPF refinement).}
  Polish the candidate to arbitrary precision using GMP \code{mpf\_t}
  arithmetic, employing a mixed Newton--Halley iteration with precision
  management modeled after Imagina.
\end{enumerate}
Three evaluation back-ends are available for Phase~A---direct iteration,
perturbation theory (PT), and linear approximation (LA)---so that the finder
can operate at any zoom depth supported by the renderer.

\subsection{Motivation: mini-Mandelbrot copies}
\label{subsec:ff-motivation}

The Mandelbrot set contains small-scale copies of itself, commonly called
\emph{mini-brots}.  Their central parameters are associated with hyperbolic
components: regions where the critical orbit
$z_0=0,\;z_{n+1}=z_n^2+c$ converges to an attracting cycle of some least
period~$p$.  The center of a given component---its \emph{nucleus}---has the
critical point on that cycle and satisfies
\begin{equation}\label{eq:ff-nucleus}
z_p(c^*) = 0,\qquad z_0 = 0,
\end{equation}
where $z_p$ denotes the $p$-fold iterate of the critical point.  Since the
recurrence is deterministic, $z_p=0$ makes the orbit repeat from its initial
state.  The equation alone gives a period dividing $p$, not necessarily least
period $p$, and can have many roots in the parameter plane.  For example,
$z_2(c)=c^2+c=c(c+1)$ has roots $0$ (least period one) and $-1$ (least
period two).  A claimed exact period must exclude returns at proper divisors
of $p$; the numerical candidate trigger is not such a certification.

Locating these nuclei is useful for several reasons:
\begin{itemize}
\item \textbf{Automatic navigation.}  Given a user click, the finder can snap
  to a candidate mini-brot center and report its period and estimated scale,
  enabling one-click deep-zoom targeting.
\item \textbf{Scale estimation.}  The derivative $\partial z_p/\partial c$ at
  the nucleus determines the local magnification scale of the mini-brot
  relative to the full set, which informs precision requirements for rendering.
\end{itemize}
At extreme zoom depths the nucleus coordinates may require thousands of bits of
precision, so the finder must combine fast low-precision search with
high-precision refinement.

\subsection{Design and implementation}
\label{subsec:ff-design}

\paragraph{Class template.}
\code{FeatureFinder<IterType, T, PExtras>} is a final class template
parameterized by the iteration-count type (\code{uint32\_t} or
\code{uint64\_t}), the working floating-point type~$T$, and a
\code{PerturbExtras} tag that selects compression behavior.  It inherits
convenience typedefs from \code{TemplateHelpers} and defines a nested
\code{Params} struct with tunable constants (maximum Newton iterations,
relative step tolerances $2^{-40}$ and $2^{-80}$, an optional absolute
residual accept threshold, and a print flag).

\paragraph{Evaluator policies (strategy pattern).}
Phase~A must iterate the critical orbit to detect the period and compute
derivatives.  Three evaluator structs encapsulate this:
\begin{itemize}
\item \code{DirectEvaluator} --- iterates $z\leftarrow z^2+c$ from scratch,
  suitable when no reference orbit is available.
\item \code{PTEvaluator} --- uses a precomputed reference orbit and
  perturbation delta $\Delta c$ to evaluate the orbit, decompressing
  reference data on the fly via \code{Runtime\-Decompressor}.
\item \code{LAEvaluator} --- uses the linear-approximation (LA) coefficients
  from \code{LAReference} to skip early iterations, falling back to
  perturbation for the remaining tail.
\end{itemize}
Each evaluator exposes a single templated method \code{Eval<FindPeriod>}.
When \code{FindPeriod=true}, it searches for the smallest $n$ at which a
periodicity trigger fires; when \code{FindPeriod=false}, it evaluates at a
known period and returns the residual, $\partial z/\partial c$, and
$z_{\mathrm{coeff}}$.

\paragraph{Common Newton loop.}
All three \code{FindPeriodicPoint} overloads delegate to
\code{FindPeriodicPoint\_Common},
which:
\begin{enumerate}
\item Calls the evaluator with \code{FindPeriod=true} to obtain a candidate
  period~$p$.
\item Applies an initial Newton correction
  $c \leftarrow c - z_p / (\partial z_p/\partial c)$ in high precision
  (the canonical $c$ is maintained as a \code{HighPrecision} value and
  projected to $T$ for each evaluation).
\item Runs up to \code{MaxNewtonIters} Newton steps in $T$-space, stopping when
  either
  $|\Delta c|^2 \le |c|^2 \cdot (2^{-40})^2$ (relative step) or
  $|z_p|^2 \le |c|^2 \cdot (2^{-40})^2$ (relative residual).
\item Performs a final correction pass and stores the result as a
  \code{PeriodicPointCandidate} inside \code{FeatureSummary}, along with the
  MPF precision needed for Phase~B.
\end{enumerate}

\paragraph{Phase~B: MPF refinement (\code{RefinePeriodicPoint\_HighPrecision}).}
The method retrieves the candidate from \code{FeatureSummary}, converts
its coordinates to GMP \code{mpf\_t} at the candidate's stored precision,
and calls the Imagina-style mixed Newton--Halley polisher described in
\cref{subsec:ff-newton,subsec:ff-halley,subsec:ff-mixed-prec}.  On
completion the refined coordinates are written back to \code{FeatureSummary}
and the candidate is marked as refined.

\paragraph{Parallel orbit evaluation.}
The high-precision forward evaluation
(\code{Evaluate\-Critical\-Orbit\-And\-Derivs}) must perform several large
\code{mpf\_mul} operations per iteration.  To overlap these, the code spawns
seven spin-based worker threads (four for derivative-precision products,
three for coordinate-precision products).  Workers spin on per-slot
\code{std::atomic\allowbreak<uint64\_t>} generation counters;
the main thread
publishes operand pointers, bumps the generation, and spin-waits on a
matching done counter.  Each slot is cache-line-aligned (\code{alignas(64)})
to avoid false sharing.

\subsection{Newton--Raphson iteration}
\label{subsec:ff-newton}

\paragraph{Root-finding formulation.}
Finding a period-$p$ nucleus amounts to solving $F(c)=0$ where
\begin{equation}\label{eq:ff-Fc}
F(c) \;=\; z_p(c),\qquad
z_0 = 0,\quad z_{n+1} = z_n^2 + c.
\end{equation}
The algorithm decomposes the root-finding problem \eqref{eq:ff-Fc} into two nested loops: an \emph{outer loop} that
updates the candidate~$c$ via Newton's method, and an \emph{inner loop} that
evaluates $F(c)$ and its derivatives by iterating the critical orbit for
$p$~steps.  The outer loop is inherently sequential---each Newton step
depends on the previous one---but the inner loop is the computational
bottleneck and is amenable to GPU acceleration
(\cref{subsec:ff-inner-gpu}).

\subsubsection{Outer loop: Newton step}
\label{subsubsec:ff-outer}

Given the outputs of the inner loop---$z_p$, $\partial z_p/\partial c$, and
optionally $\partial^2 z_p/\partial c^2$---Newton's method produces the
update
\begin{equation}\label{eq:ff-newton-step}
\Delta c \;=\; \frac{F(c)}{F'(c)}
         \;=\; \frac{z_p}{\partial z_p / \partial c},
\qquad c \leftarrow c - \Delta c.
\end{equation}
The first-order Taylor model is
$F(c-\Delta c)\approx F(c)-F'(c)\Delta c$.  Setting that model to zero
gives the correction in \eqref{eq:ff-newton-step}; another complete orbit
evaluation then checks the new candidate rather than assuming the linear
model was exact.
Writing $z_p = z_x + i\,z_y$ and $\partial z_p/\partial c = d_x + i\,d_y$,
the complex division expands to:
\begin{align}
\Delta c_x &= \frac{z_x\,d_x + z_y\,d_y}{d_x^2 + d_y^2}, \\[4pt]
\Delta c_y &= \frac{z_y\,d_x - z_x\,d_y}{d_x^2 + d_y^2}.
\end{align}
(When the Halley gate is satisfied, a higher-order step using
$\partial^2 z_p/\partial c^2$ replaces the Newton step; see
\cref{subsec:ff-halley}.)

\paragraph{A worked period-two Newton step.}
For $p=2$, the inner loop gives
$F(c)=c^2+c$, $F'(c)=2c+1$, and $F''(c)=2$.
Starting from $c=-0.9$ gives $F=-0.09$ and $F'=-0.8$.
The correction is $\Delta c=0.1125$, so the new parameter is
$-1.0125$.  Two more Newton steps give approximately
$-1.00015243902439$ and $-1.00000002323057$.
Writing $c=-1+e$ shows exactly why the errors decrease quadratically:
\[
e_{\mathrm{next}}
=e-\frac{e(e-1)}{2e-1}
=\frac{e^2}{2e-1}.
\]
Near $e=0$, the denominator is close to $-1$ while the numerator is
quadratic.  At $c=-1/2$, however, $F'=0$ and the Newton division is
undefined; near a simple root is an essential assumption.
A starting point near $c=0$ can instead converge to the lower-period root.

\paragraph{Convergence.}
Standard Newton--Raphson converges quadratically near a simple root.  The
Phase~A loop in \code{FindPeriodicPoint\_Common} uses two independent
stopping criteria, whichever triggers first:
\begin{enumerate}
\item \emph{Relative step:}
  $|\Delta c|^2 \le |c|^2 \cdot \epsilon_{\mathrm{step}}^2$
  with $\epsilon_{\mathrm{step}} = 2^{-40}$.
\item \emph{Relative residual:}
  $|z_p|^2 \le |c|^2 \cdot \epsilon_{\mathrm{diff}}^2$
  with $\epsilon_{\mathrm{diff}} = 2^{-40}$.
\end{enumerate}
An optional absolute residual accept (\code{Eps2Accept}) provides a third
early-out that can be enabled via \code{Params}.

\paragraph{Period detection trigger.}
Before Newton iteration begins, the evaluator must identify the candidate
period.  The criterion used is: at iteration~$n$, if
\begin{equation}\label{eq:ff-period-trigger}
|z_n|^2 \;<\; R^2 \cdot \left|\frac{\partial z_n}{\partial c}\right|^2,
\end{equation}
then $p = n$ is accepted as the candidate period, where $R$ is the search
radius.  Write $d_n=\partial z_n/\partial c$ for the derivative at the
current scan index.  The inequality is equivalent to $|z_n/d_n|<R$ when
$d_n\ne0$: the linearized Newton displacement for a return at iteration $n$
fits inside the search radius.  It proposes a candidate rather than proving that a root lies there.
The test detects the moment the orbit passes close to the origin
relative to the local stretching factor $|\partial z/\partial c|$, which is
characteristic of near-periodic behavior.  A dynamic tightening mechanism
(implemented in the \code{PeriodicityPP} helper) progressively shrinks the
effective $R^2$ during the scan to reduce false triggers.

\subsubsection{Inner loop: forward evaluation}
\label{subsubsec:ff-inner}

\begin{figure}[!htbp]
\centering
\begin{tikzpicture}[
    >=Stealth,
    box/.style={draw, rounded corners, minimum width=2.8cm, minimum height=0.7cm,
                align=center, font=\small},
    outerbox/.style={box, fill=blue!8},
    innerbox/.style={box, fill=orange!10},
    precbox/.style={font=\scriptsize\itshape, text=gray},
    every edge/.style={draw, ->},
]
  % Outer loop
  \node[outerbox] (init) {Initial guess $c$};
  \node[outerbox, below=0.5cm of init] (inner) {Inner loop\\(evaluate $z_p$, $\mathit{dzdc}_p$, $d2_p$)};
  \node[outerbox, below=0.5cm of inner] (step) {Newton step\\$c \leftarrow c - z_p/\mathit{dzdc}_p$};
  \node[outerbox, below=0.5cm of step] (conv) {Converged?};
  \node[outerbox, below=0.5cm of conv] (done) {Return $c^*$};

  \draw[->] (init) -- (inner);
  \draw[->] (inner) -- (step);
  \draw[->] (step) -- (conv);
  \draw[->] (conv) -- node[right, font=\scriptsize] {yes} (done);
  \draw[->] (conv.west) -- ++(-1.5,0) |- node[near start, left, font=\scriptsize] {no} (inner.west);

  % Inner loop detail
  \node[innerbox, right=2.5cm of inner, yshift=1.0cm] (orbit) {$z \leftarrow z^2 + c$};
  \node[precbox, right=0.1cm of orbit.east] {coord};
  \node[innerbox, below=0.3cm of orbit] (dzdc) {$\mathit{dzdc} \leftarrow 2z \cdot \mathit{dzdc} + 1$};
  \node[precbox, right=0.1cm of dzdc.east] {deriv};
  \node[innerbox, below=0.3cm of dzdc] (d2) {$d2 \leftarrow 2(\mathit{dzdc}^2 + z \cdot d2)$};
  \node[precbox, right=0.1cm of d2.east] {${\sim}53$-bit};

  \node[font=\scriptsize, above=0.15cm of orbit] (repeat) {Repeat $p$ times:};

  \begin{scope}[on background layer]
    \node[draw, dashed, rounded corners, fill=orange!4,
          fit=(orbit)(dzdc)(d2)(repeat), inner xsep=6pt, inner ysep=10pt,
          label={[font=\scriptsize\bfseries]above:Inner loop (GPU-amenable)}] (innerfit) {};
  \end{scope}

  \draw[->, dashed] (inner.east) -- (innerfit.west);

  % Outer loop label
  \begin{scope}[on background layer]
    \node[draw, dotted, rounded corners, fill=blue!3,
          fit=(init)(done)(conv), inner xsep=8pt, inner ysep=8pt,
          label={[font=\scriptsize\bfseries]above:Outer loop (CPU)}] {};
  \end{scope}
\end{tikzpicture}
\caption[Newton--Raphson two-level loop structure]{Newton--Raphson two-level loop structure.  The outer loop
  (left) updates the candidate~$c$ via Newton steps on the CPU.  Each
  step calls the inner loop (right), which iterates $p$~times to
  evaluate $z_p$, $\mathit{dzdc}_p$, and $d2_p$ at three different
  precision levels. All right-hand sides use the previous iteration's state.
  The inner loop is the computational bottleneck and
  the target for GPU acceleration.}
\label{fig:newton-loops}
\end{figure}

Each Newton step requires evaluating $F(c) = z_p(c)$ and its first two
derivatives at the current guess~$c$, as shown in \cref{fig:newton-loops}.
This evaluation is done by iterating the
critical orbit from $z_0 = 0$ for exactly $p$~steps, accumulating three
quantities simultaneously at \emph{different precision levels}:

\begin{enumerate}
\item \textbf{Orbit} ($z$) --- \emph{coordinate precision} (\code{mpf\_t},
  hundreds to tens of thousands of bits):
  $z_{n+1} = z_n^2 + c$.
\item \textbf{First derivative} ($\partial z/\partial c$, written
  $\mathit{dzdc}$) --- \emph{derivative precision} (\code{mpf\_t}, a
  separately selected precision $P_d$, typically reduced relative to $P_c$
  but not always bounded by it; see \cref{subsec:ff-mixed-prec}):
  $\mathit{dzdc}_{n+1} = 2\,z_n \cdot \mathit{dzdc}_n + 1$.
\item \textbf{Second derivative} ($\partial^2 z/\partial c^2$, written
  $d2$) --- \emph{low precision} (\code{HDRFloat<double>}, ${\sim}53$-bit
  mantissa with wide exponent range):
  $d2_{n+1} = 2\bigl(\mathit{dzdc}_n^{\,2} + z_n \cdot d2_n\bigr)$.
\end{enumerate}

\noindent
Initial conditions: $z_0 = 0$, $\;\mathit{dzdc}_0 = 0$, $\;d2_0 = 0$.
Differentiating $z_{n+1}=z_n^2+c$ once gives
$d_{n+1}=2z_nd_n+1$.  Differentiating again and applying the product rule
gives $q_{n+1}=2(d_n^2+z_nq_n)$, where $q_n=\partial^2z_n/\partial c^2$.
For period two, the states are $(z_1,d_1,q_1)=(c,1,0)$ and
$(z_2,d_2,q_2)=(c^2+c,2c+1,2)$, recovering the derivatives in the worked
Newton example without differentiating a large degree-$2^{p-1}$ polynomial.

The second derivative is deliberately kept at low precision because it
enters the Halley correction factor (\cref{subsec:ff-halley}) and a local
error proxy (\cref{eq:ff-err}), rather than the coordinate recurrence.
The gate and stop decision compare exponents of quantities formed from it;
the Halley correction itself also depends on its mantissa and phase.
The usefulness of a 24- or 53-bit approximation therefore depends on the
conditioning of the derivative recurrence, as analyzed in
\cref{subsec:ff-halley}.
The $z$ and $\mathit{dzdc}$ values are truncated to low-precision
\code{HDRFloat} before entering the $d2$ recurrence---the $d2$
computation never touches the high-precision \code{mpf\_t}
representation.

After $p$~iterations the outputs $z_p$, $\mathit{dzdc}_p$, and $d2_p$ are
returned to the outer loop.

\paragraph{Real/imaginary expansion.}
To implement the inner loop with ordinary (real-valued) arithmetic, write
\[
z = x + i\,y,\qquad
\mathit{dzdc} = dx + i\,dy,\qquad
d2 = d2x + i\,d2y,\qquad
c = c_x + i\,c_y.
\]
Each iteration step then consists of the following scalar operations.

\medskip\noindent
\textbf{Orbit}\quad ($z \leftarrow z^2 + c$):
\begin{equation}\label{eq:ff-orbit-ri}
\boxed{\begin{aligned}
x_{n+1} &= x_n^2 - y_n^2 + c_x, \\
y_{n+1} &= 2\,x_n\,y_n + c_y.
\end{aligned}}
\end{equation}

\noindent
\textbf{First derivative}\quad ($\mathit{dzdc} \leftarrow 2\,z\cdot\mathit{dzdc} + 1$):
\begin{equation}\label{eq:ff-dzdc-ri}
\boxed{\begin{aligned}
dx_{n+1} &= 2\,x_n \cdot dx_n - 2\,y_n \cdot dy_n + 1, \\
dy_{n+1} &= 2\,x_n \cdot dy_n + 2\,y_n \cdot dx_n.
\end{aligned}}
\end{equation}

\noindent
\textbf{Second derivative}\quad
($d2 \leftarrow 2\,(\mathit{dzdc}^{\,2} + z \cdot d2)$):
\begin{equation}\label{eq:ff-d2-ri}
\boxed{\begin{aligned}
d2x_{n+1} &= 2\,dx_n^2 - 2\,dy_n^2
              + 2\,x_n \cdot d2x_n - 2\,y_n \cdot d2y_n, \\
d2y_{n+1} &= 4\,dx_n \cdot dy_n
              + 2\,x_n \cdot d2y_n + 2\,y_n \cdot d2x_n.
\end{aligned}}
\end{equation}

\noindent
Each line is a sequence of multiplies, adds, and subtracts on real-valued
scalars---no complex arithmetic library is required.  Note that
\emph{within} a single iteration, the orbit update uses $x_n$,~$y_n$ (the
old values), the $\mathit{dzdc}$ update uses $x_n$,~$y_n$ and the old
$dx_n$,~$dy_n$, and the $d2$ update uses the old $x_n$,~$y_n$,
$dx_n$,~$dy_n$, $d2x_n$,~$d2y_n$.  The $d2$ recurrence must therefore be
evaluated \emph{before} the orbit and $\mathit{dzdc}$ values are
overwritten (or all old values must be saved).

\paragraph{Operation count.}
Not all operations are equally expensive.  The CPU orbit and first-derivative
recurrences use high-precision \code{mpf\_t} arithmetic, where large
$N$-limb multiplication costs substantially more than the $O(N)$
add/subtract operations.  MPIR chooses multiplication algorithms according
to operand size; a single asymptotic formula does not describe every precision.
The second derivative uses low-precision
\code{HDRFloat} arithmetic (${\sim}53$-bit on the CPU path,
${\sim}24$-bit on the GPU path)---a handful of scalar floating-point
operations per iteration, negligible by comparison.

Counting real-valued scalar operations in the expanded CPU schedule from
\cref{eq:ff-orbit-ri,eq:ff-dzdc-ri,eq:ff-d2-ri}, each iteration requires:

\begin{center}
\small
\begin{tabular}{llccc}
\hline
\textbf{Recurrence} & \textbf{Precision} &
  \textbf{Muls} & $\boldsymbol{\times 2}$ & \textbf{Add/subs} \\
\hline
Orbit ($z^2 + c$) & coord &
  3 & 1 & 3 \\
$\mathit{dzdc}$ ($2z \cdot \mathit{dzdc} + 1$) & deriv &
  4 & 2 & 3 \\
$d2$ ($2(\mathit{dzdc}^2 {+} z {\cdot} d2)$) & low-prec HDR &
  7 & 3 & 5 \\
\hline
\textbf{HP total} & &
  \textbf{7} & \textbf{3} & \textbf{6} \\
\hline
\end{tabular}
\end{center}

\noindent
``Muls'' in the table are full CPU $N \times N$ multiplications
(\code{mpf\_mul})---the dominant cost at large precision.
``$\times 2$'' denotes an $O(N)$ scaling by~two, and ``Add/subs'' includes
constant additions.  The ``HP total'' row counts only the high-precision
coordinate and first-derivative work; the $d2$ row contributes only cheap
low-precision operations.

The multithreaded CPU evaluator dispatches the seven expanded products to
seven spin-based worker threads (\cref{subsec:ff-threading}).  The fused v2 GPU
engine uses different algebra.  It forms the orbit with
$(x+y)(x-y)$ and $2xy$, then evaluates the complex first-derivative product
with the three-product Gauss form.  It therefore performs five pointwise
Montgomery products after transforming the shared input streams, rather than
launching seven standalone high-precision multiplications.  Factors of two,
$c_x$, $c_y$, and the derivative's $+1$ enter as aligned linear terms in the
fused output construction (\cref{subsec:ff-inner-gpu}).

The $d2$ computation runs on the CPU main thread or one GPU thread using
truncated copies of~$z$ and $\mathit{dzdc}$.  Its low-precision work completes
alongside the expensive high-precision recurrence.

\paragraph{The \code{zcoeff} product.}
In addition to $\partial z_p/\partial c$, the code tracks a multiplicative
coefficient
\begin{equation}\label{eq:ff-zcoeff}
w_n = \prod_{k=1}^{n-1} 2\,z_k,
\end{equation}
with the convention $w_0 = w_1 = 1$.  Expanding
$d_{n+1}=2z_nd_n+1$ repeatedly shows that $w_n$ is the contribution
propagated from $d_1=1$; subsequent $+1$ terms also contribute to $d_n$.
For instance, $d_3=4z_2z_1+2z_2+1=w_3+2z_2+1$.
The product is used later
to compute the intrinsic scale of the mini-brot (see
\cref{subsec:ff-intrinsic-radius}).  In the code, \code{zcoeff} is
initialized to~$1$ at $n=0$ (because $z_0 = 0$ would otherwise collapse the
product) and then multiplied by $2z$ at each subsequent step, matching the
recurrence for the leading factor of the derivative.

\subsubsection{GPU acceleration of the inner loop}
\label{subsec:ff-inner-gpu}

The fused v2 backend executes the orbit and first-derivative recurrences in the
same persistent cooperative kernel.  Orbit-only work uses TwoWay mode for
$x$ and~$y$; Newton--Raphson work uses FourWay mode for $x$, $y$, $dx$, and
$dy$.  The mathematical product reduction, NTT ordering, CUDA memory hierarchy,
shared barrier schedule, signed carry, and normalization are developed in
\cref{subsec:fused-ntt-reference-step,fig:ntt-pipeline}.

For Feature Finder evaluation, FourWay mode adds the three Gauss-form products
in~\eqref{eq:fused-derivative-products} to the two orbit products.  It reconstructs
the first derivative with~\eqref{eq:fused-derivative-reconstruction}; the aligned
$+1$ enters during signed gathering rather than through a separate transform or
addition kernel.  Enabling the derivative streams widens the work performed
between the existing cooperative-grid boundaries but does not add an independent
high-precision pipeline.

A leader thread separately computes the second-derivative update
(\cref{eq:ff-d2-ri}) in approximately 24-bit \code{HDRFloat<float>} arithmetic from
truncated copies of the current $z$ and $\mathit{dzdc}$.  This inexpensive
update precedes the fused high-precision phases and does not introduce another
cooperative-grid dependency.

\medskip
\noindent
The outer Newton loop (typically ${\sim}10$--$30$ iterations) runs on
the host; the inner loop---$p$ iterations of high-precision
arithmetic per Newton step---runs on the GPU.  The host reads back
$z_p$, $\mathit{dzdc}_p$, and $d2_p$ after each $p$-iteration batch,
computes the Newton/Halley step, updates~$c$, and relaunches the inner
loop.  For large $p$ and $N$, the transform arithmetic volume is approximately
$O(p \cdot N \log N)$ per Newton step when aligned support scales with
limb count, motivating GPU acceleration at suitable precisions.  \code{Orbit\-Endpoint\-Evaluator} maps the coordinate
precision to a supported limb count with \code{BitsToSupportedLimbCount},
selects the corresponding Newton--Raphson precision specialization through
\code{DispatchByLimbCount}, and calls
\code{EvaluateCriticalOrbitAndDerivs\_GPU} to run the fused kernel.

\subsection{Halley's method}
\label{subsec:ff-halley}

The second derivative $F'' = \partial^2 z_p/\partial c^2$ is computed during
every forward evaluation, even when the Halley step itself is not taken.  This
is because $F''$ serves two distinct purposes in the refinement loop:

\begin{enumerate}
\item \textbf{Halley correction.}  When the Halley gate (\cref{eq:ff-halley-gate})
  is satisfied, $F''$ enables a cubic-convergence step that can save expensive
  forward evaluations compared to plain Newton (see below).
\item \textbf{Error estimator.}  The stopping criterion
  (\cref{eq:ff-err}) uses $|F''|^2$ to estimate the local coordinate error
  after a step.  It avoids an extra complete orbit evaluation solely to
  obtain a new residual; the estimate remains a numerical proxy.
\end{enumerate}

Because $F''$ is carried in \code{HDRFloat} (${\sim}53$~bits of
mantissa on the CPU path, ${\sim}24$~bits on the GPU path; both with a
wide exponent), it adds only cheap low-precision arithmetic to
the forward evaluation: a few \code{HDRFloat} multiplies and additions per
iteration of the period loop, compared with the expensive \code{mpf\_t}
multiplies needed for $z$ and $\partial z/\partial c$.  The cost is therefore
negligible relative to the rest of the evaluation, making it worthwhile to
compute unconditionally.

Phase~B optionally upgrades from Newton to Halley iteration for cubic
convergence.  The Halley step for $F(c) = z_p(c)$ is
\begin{equation}\label{eq:ff-halley-step}
\Delta_H \;=\;
  \frac{2\,F(c)\,F'(c)}
       {2\bigl(F'(c)\bigr)^2 - F(c)\,F''(c)},
\qquad c \leftarrow c - \Delta_H,
\end{equation}
where $F'=\partial z_p/\partial c$ and $F''=\partial^2 z_p/\partial c^2$.
A quadratic Taylor model gives
$0\approx F-F'\Delta c+\tfrac12F''(\Delta c)^2$.
Rearranging it and replacing one factor of $\Delta c$ in the quadratic term
by the Newton correction $\Delta_N=F/F'$ gives
\[
\Delta c\approx\Delta_N+\frac{F''}{2F'}\Delta_N\Delta c,
\qquad
\Delta c\approx\frac{\Delta_N}{1-FF''/[2(F')^2]}.
\]
The resulting rational expression is precisely the Halley update above.
The construction incorporates curvature without requiring a separate
quadratic-root choice for a complex correction.

To evaluate \eqref{eq:ff-halley-step}, the denominator $D = 2(F')^2 - F\,F''$ involves a complex squaring and a
complex product.  Writing $F = f_x + i\,f_y$, $F' = f'_x + i\,f'_y$, and
$F'' = f''_x + i\,f''_y$:
\begin{align}
D_x &= 2\,(f'^2_x - f'^2_y) - (f_x\,f''_x - f_y\,f''_y), \\
D_y &= 2\,(2\,f'_x\,f'_y) - (f_x\,f''_y + f_y\,f''_x).
\end{align}
The numerator $N = 2\,F\,F'$ is a complex product, and the final step
$\Delta_H = N / D$ is a complex division as in \cref{eq:ff-newton-step}.

\paragraph{Second derivative recurrence.}
The second derivative $d2 = \partial^2 z/\partial c^2$ is computed by the
inner loop (\cref{subsubsec:ff-inner}) using the recurrence in
\cref{eq:ff-d2-ri}.  In the code this quantity is carried as a pair of
\code{HDRFloat} values (real and imaginary parts), which provide a
wide exponent range at low mantissa precision---\code{HDRFloat<double>}
(${\sim}53$-bit) on the CPU path and \code{HDRFloat<float>}
(${\sim}24$-bit) on the GPU path.  The coordinate-precision
quantities $z$ and $\partial z/\partial c$ are stored in \code{mpf\_t},
while $\partial^2 z/\partial c^2$ remains in HDR throughout the forward
evaluation.  When the Halley step is computed, the HDR second
derivative is promoted to \code{mpf\_t} at coordinate precision.

\paragraph{Halley gate.}
Halley's step is only beneficial when $F''$ is well-conditioned relative to
$F$ and $F'$.  The code uses the dimensionless ratio
\begin{equation}\label{eq:ff-halley-gate}
\rho^2 \;=\;
  \frac{|F|^2\,|F''|^2}{|F'|^4}
\end{equation}
and activates Halley when the reduced HDR exponent
$E_{\rho^2}\le-12$.  For a positive normalized value
$\rho^2=m2^{E_{\rho^2}}$ with $1\le m<2$, this means
$\rho^2<2^{-11}$; it is an exponent-bin gate rather than a literal
comparison with $2^{-12}$.  When $\rho^2$ is
small the denominator $2(F')^2 - F\,F''$ is dominated by the $(F')^2$ term,
so the Halley correction is a small, well-determined perturbation of the
Newton step.  If the Halley denominator turns out to be exactly zero (or the
gate is not satisfied), the code falls back to a plain Newton step for that
iteration.

\paragraph{A period-two Halley comparison.}
At $c=-0.999$, $F=-0.000999$, $F'=-0.998$, and $F''=2$.
Here $\rho^2\approx4.02\times10^{-6}$, well inside the gate.
The Newton update is approximately $-1.00000100200401$.
Halley instead gives $-0.999999998996994$.
Indeed, for $c=-1+e$ its exact error is
\[
e_{\mathrm{next,Halley}}=\frac{e^3}{1-3e+3e^2}.
\]
At $e=0.001$, this is approximately $1.00300601\times10^{-9}$,
whereas Newton leaves error approximately $-1.00200401\times10^{-6}$.
The extra derivative can therefore save a full period-length evaluation.

\paragraph{Is Halley worthwhile given low-precision $F''$?}
The second derivative $F''$ is accumulated in \code{HDRFloat}
throughout the forward evaluation, providing only $b$~bits of
mantissa (with a wide exponent range): $b \approx 53$ on the CPU path
(\code{HDRFloat<double>}) and $b \approx 24$ on the GPU path
(\code{HDRFloat<float>}, since \code{SharkParamsNR} uses
\code{SubType=float}).  Meanwhile $F$ and $F'$ are carried in
\code{mpf\_t} at hundreds or thousands of bits.  This precision mismatch
limits the benefit of the Halley correction.

The Halley step can be decomposed as
\begin{equation}\label{eq:ff-halley-decomp}
\Delta_H \;=\; \Delta_N \cdot \frac{1}{1 - \frac{F \cdot F''}{2(F')^2}}
\;=\; \Delta_N \cdot \frac{1}{1 - \tfrac{1}{2}\rho\, e^{i\theta}},
\end{equation}
where $\Delta_N = F/F'$ is the Newton step, $\rho = \sqrt{\rho^2} =
|F|\cdot|F''|/|F'|^2$ is the first-power ratio (the gate criterion checks
$\rho^2$; the correction factor has magnitude $\rho/2$), and $e^{i\theta}$
is the complex phase of $F \cdot F'' / (F')^2$.

\paragraph{Why it works: error analysis of the correction factor.}
Starting from \eqref{eq:ff-halley-decomp}, suppose the computed second
derivative has a small relative error,
$\hat{F}'' = F''(1 + \epsilon)$.  A well-conditioned calculation can make
$|\epsilon|$ comparable to $2^{-b}$, but $b$ mantissa bits alone do not
prove that bound after a long recurrence with cancellation.  Also assume
$F,F'$ are sufficiently accurate and the root is simple.  The
Halley step with the approximate $\hat{F}''$ is
\[
\hat{\Delta}_H \;=\; \Delta_N \cdot
  \frac{1}{1 - \frac{F \cdot \hat{F}''}{2(F')^2}}.
\]
Near a simple root, the error $e_k = c_k - c^*$ satisfies $F \approx F' e_k$,
so the correction ratio
\[
\frac{F \cdot F''}{2(F')^2} \;\approx\;
  \frac{e_k \cdot F''}{2 F'} \;=\; O(e_k).
\]
This term shrinks as $c_k$ converges.  The \emph{error} introduced by using
$\hat{F}''$ instead of $F''$ modifies this ratio by
\[
\frac{F \cdot \hat{F}''}{2(F')^2}
  - \frac{F \cdot F''}{2(F')^2}
  \;=\; \frac{F \cdot F'' \epsilon}{2(F')^2}
  \;=\; O(e_k \cdot \epsilon).
\]
This quantity is the product of the current error and the $b$-bit truncation noise.
The resulting Halley step error, relative to the true Halley step, is therefore
$O(e_k \cdot \epsilon)$, producing a next error of
\[
e_{k+1}^{\text{(approx Halley)}} \;=\;
  \underbrace{O(e_k^3)}_{\text{true Halley}}
  \;+\;
  \underbrace{O(e_k^2 \cdot \epsilon)}_{\text{truncation of } F''}.
\]
(The extra factor of $e_k$ in the second term comes from the fact that
$\Delta_N$ itself is $O(e_k)$, and the relative perturbation of the
correction factor is $O(e_k \cdot \epsilon)$.)

The constants in these order estimates depend on the local derivatives.
The meaningful error scale is normalized: with
$\kappa=|F''(c^*)/F'(c^*)|$ and $\eta_k=\kappa|e_k|$, assuming
$\kappa\ne0$ and moderate higher-derivative ratios, a local model gives
$\eta_{k+1}=O(\eta_k^3)+O(\eta_k^2|\epsilon|)$.
Here $\kappa$ converts coordinate error into a dimensionless measure of
nonlinearity.  Under the further model $|\epsilon|\sim2^{-b}$, two regimes emerge:
\begin{itemize}
\item \textbf{Early convergence} ($\eta_k \gg 2^{-b}$): the $O(e_k^3)$ term
  dominates, giving the same cubic order as full-precision Halley; the
  second-derivative noise is smaller than the cubic term in this model.
\item \textbf{Late convergence} ($\eta_k \ll 2^{-b}$): the
  $O(e_k^2 \cdot \epsilon)$ term dominates.  Since $\epsilon$ is a fixed
  ${\sim}2^{-b}$, the effective recurrence becomes $|e_{k+1}| \approx
  C \cdot |e_k|^2 \cdot 2^{-b}$, which is quadratic convergence (like Newton)
  but with an extra factor of $2^{-b}$ in this model.  In terms of normalized bits,
  this model gives approximately $2n+b$ normalized bits per step---still better than Newton's
  ${\sim}2n$, but no longer cubic.
\end{itemize}

In normalized units, the crossover occurs near
$\eta_k\approx2^{-b}$, or $|e_k|\approx2^{-b}/\kappa$.
Thus ${\sim}b$ \emph{normalized} bits mark the crossover in this model;
the coordinate error may already have approximately $b+\log_2\kappa$
correct bits.  Strong cancellation or inaccurate $F'$ can change the model.

\paragraph{Empirical confirmation ($b = 53$).}
Numerical experiments on the period-3 nucleus illustrate this local model
for the CPU path ($b = 53$, \code{HDRFloat<double>}):

\begin{center}
\begin{tabular}{crrr}
\hline
\textbf{Step} & \textbf{Newton} & \textbf{Halley (53-bit $F''$)}
              & \textbf{Halley (full $F''$)} \\
\hline
0 &  13.8 &  20.7 &  20.7 \\
1 &  26.5 &  60.4 &  60.4 \\
2 &  51.9 & 173.8 & 179.6 \\
3 & 102.6 & 400.6 & 536.9 \\
4 & 204.1 & 854.2 & 1608.9 \\
5 & 407.0 & 1761.4 & --- \\
6 & 812.8 & --- & --- \\
7 & 1624.4 & --- & --- \\
\hline
\end{tabular}
\end{center}

(Correct bits of the period-3 nucleus, starting from the same initial guess.)
Steps~0--1 show cubic convergence for both Halley variants (ratio
${\approx}\,3$). At step~2, the 53-bit Halley begins to diverge from the
full-precision version (173.8 vs.\ 179.6), following its crossover near
53 normalized bits.  By step~4, the 53-bit Halley's ratio has dropped to
${\sim}2.1$, approaching quadratic.

\paragraph{Net benefit.}
On the period-3 example, the low-precision Halley reaches 1761~correct bits
in 6~steps versus Newton's 8~steps for 1624~bits---saving two full forward
evaluations.  This saving comes from the early cubic burst: by the time the
method degenerates to quadratic, it already has a ${\sim}3\times$ lead in
correct bits over Newton, and this lead persists through the later quadratic
phase.  Such savings are most pronounced when Phase~B starts below the
${\sim}b$-normalized-bit crossover.  At already-accurate starts, the large cubic lead is
absent, so the retained benefit is narrower: representative runs show only a
0--1 outer-iteration advantage in this late regime, while $F''$ still remains
useful for the Halley gate and error estimator.  For large period~$p$, even
that one saved forward evaluation still avoids $p$~expensive
high-precision multiplies.

\paragraph{Dependence on initial-guess quality.}
The important starting-point measure is error relative to the local feature
scale, not zoom factor or coordinate bits alone.  Phase~A supplies a numerical
candidate, and the local ratio $\kappa=|F''/F'|$ can grow greatly with
period.  If $\kappa\approx2^{88}$ and the absolute coordinate error is
$2^{-100}$, the normalized error is approximately $2^{-12}$.
With a well-conditioned 53-bit second derivative, that point is still in
the cubic regime of the local model.  An error of $2^{-150}$ instead gives
normalized error $2^{-62}$ and lies beyond its crossover.  Pixel spacing,
candidate selection, derivative conditioning, and Phase~A refinement all
affect the starting error; zoom alone does not determine the iteration saving.

\paragraph{Interpreting the diagnostic metric.}
The diagnostic ratio $\rho^2$ (\cref{eq:ff-halley-gate}) measures
convergence in a normalized sense rather than counting correct bits of
the candidate~$c$.  Near a simple root, $F \approx F' e_k$ where
$e_k = c_k - c^*$, so
\begin{equation}\label{eq:ff-rho2-offset}
\rho^2 \;\approx\; |e_k|^2 \cdot \frac{|F''|^2}{|F'|^2},
\qquad
  \text{normalized bits}
  \;=\; -\tfrac{1}{2}\log_2 \rho^2
  \;\approx\; \text{correct bits} - \log_2 \frac{|F''|}{|F'|}.
\end{equation}
For the period-3 example, $|F''/F'| \approx 2^{2.2}$, giving a small
2.2-bit offset.  At period~16045, $|F''/F'| \approx 2^{88}$, so the
$\rho^2$-derived metric understates actual correct bits by approximately
88~bits in that example.  The normalized metric measures the local
nonlinearity that controls the model crossover; it is not a count of correct
coordinate bits.  The implementation reports an exponent-based integer
approximation to that metric.

\paragraph{GPU variant ($b = 24$).}
The GPU NR path (\code{SharkParamsNR}) uses \code{SubType=float},
so $d2$ is accumulated in \code{HDRFloat<float>} with approximately
24 mantissa bits.  Under the same well-conditioned local model, its crossover
occurs near 24 normalized bits rather than 53.  Retaining more second-derivative
precision could extend the cubic regime, but would also increase arithmetic
latency and register use in the leader thread.  The current implementation
chooses the cheaper representation; a universal outer-step saving does not
follow from its mantissa width.

The wide exponent range prevents ordinary native underflow and overflow in
these auxiliary values.  It does not make their exponent mathematically exact:
rounding and cancellation can change the computed magnitude, including its
exponent bin.  The Halley gate and error proxy therefore remain numerical
decisions, even though their final comparisons use exponents.

\subsection{Mixed-precision refinement}
\label{subsec:ff-mixed-prec}

The Phase~B polisher (\code{RefinePeriodicPoint}) operates with GMP
\code{mpf\_t} coordinate arithmetic on the host and a selectable CPU or GPU
endpoint evaluator.  It uses two distinct high-precision levels to balance cost
and accuracy, following the strategy used in Imagina.

\paragraph{Coordinate vs.\ derivative precision.}
The candidate coordinate $c$ and the orbit value $z_p$ are maintained at full
\emph{coordinate precision} $P_c$ (typically hundreds to tens of thousands of
bits, determined by the zoom depth).  The derivative
$\partial z_p/\partial c$ is maintained at a reduced \emph{derivative
precision} $P_d$ chosen by
\begin{equation}\label{eq:ff-deriv-prec}
P_d = \operatorname{clamp}\!\left(
  \frac{-E_{\mathrm{scale}} + 32}{4},\;
  256,\;
  P_c + E_{\max}^{|c|} + 32
\right),
\end{equation}
where $E_{\mathrm{scale}}$ is the base-2 exponent of the scale factor
$\text{Scale} \approx 1/|w_p \cdot \partial z_p/\partial c|$ and
$E_{\max}^{|c|}$ is the exponent of $\max(|\Re c|,|\Im c|)$.
The derivative must be accurate enough for the desired coordinate correction.
Reducing its precision can lower the cost of four CPU products per orbit step,
but its rounding error can also limit convergence.  The clamped rule is an
implementation policy that balances those costs rather than a proof of the
required derivative accuracy.  Integer precision calculations round the
intermediate division before clamping.

\paragraph{Error estimator and stopping criterion.}
After each step $\Delta c$, the code computes an error proxy entirely in
\code{HDRFloat<double>}:
\begin{equation}\label{eq:ff-err}
\mathrm{err} \;=\;
  \frac{|\Delta c|^4 \cdot |F''|^2}{|F'|^2}.
\end{equation}
Near a simple root, Taylor expansion gives the local Newton estimate
$|e_{k+1}| \approx |F''|/(2|F'|) \cdot |\Delta c|^2$,
so $\mathrm{err} \approx 4\,|e_{k+1}|^2$ when the derivatives and step
are accurate.  The relation is an asymptotic estimate, not a classical
certified bound.  For the period-two Newton step at $c=-0.999$,
$|\Delta c|\approx0.001001002$ and $|F''/F'|\approx2.004008$,
giving $\mathrm{err}\approx4.032\times10^{-12}$; four times the actual
next squared error is approximately $4.016\times10^{-12}$.
Iteration stops when
\begin{equation}\label{eq:ff-stop}
-E_{\mathrm{err}} \;\ge\; 2\,P_c,
\end{equation}
where $E_{\mathrm{err}}$ is the reduced HDR exponent.  For nonzero
$\mathrm{err}=m2^{E_{\mathrm{err}}}$ with $1\le m<2$, the selected bin
implies $\mathrm{err}<2^{1-2P_c}$.  The intended coordinate error scale
is therefore approximately $2^{-P_c}$ under the local model.  Reduced
second-derivative arithmetic, Halley steps, and finite-precision rounding
prevent this stop decision from certifying $P_c$ correct bits.

\paragraph{Final correction and accept/reject.}
After the main loop converges (or exhausts the iteration budget), a single
additional Newton step is applied as a final correction.  The refined $c$ is
then checked against the initial seed $c_0$: if
$|c - c_0|^2 > R^2$ (the squared search radius, carried in \code{mpf\_t}),
the result is rejected and $c$ is reset to~$c_0$.  This guard prevents the
iteration from wandering to a distant nucleus outside the user's region of
interest.

\subsection{Intrinsic radius and scale}
\label{subsec:ff-intrinsic-radius}

Once the nucleus $c^*$ and its period~$p$ are known, the local magnification
of the mini-brot copy is characterized by
\begin{equation}\label{eq:ff-scale}
\mathrm{Scale} \;=\;
  \frac{1}{\bigl|w_p \cdot \partial z_p/\partial c\bigr|},
\end{equation}
where $w_p = \prod_{k=1}^{p-1} 2\,z_k$ is the \code{zcoeff} product
\eqref{eq:ff-zcoeff} evaluated at $c^*$.  The local return map explains the
inverse product: for a small initial state $u$ and parameter change
$\delta c$, it has the leading form
$u'\approx w_pu^2+d_p\delta c$, where $d_p=\partial z_p/\partial c$.
The initial square produces $u^2$ and the subsequent derivatives multiply it
by $w_p$.  Rescaling to $v=w_pu$ gives
$v'\approx v^2+(w_pd_p)\delta c$, so a unit change of the rescaled
parameter corresponds to $1/|w_pd_p|$ in the original plane.  Higher-order
terms remain omitted in this local model.
The code reports an \emph{intrinsic radius}
\begin{equation}\label{eq:ff-intrinsic-radius}
r_{\mathrm{int}} = 4 \cdot \mathrm{Scale}
  = \frac{4}{\bigl|w_p \cdot \partial z_p/\partial c\bigr|}.
\end{equation}
The quantity is a renderer navigation scale, despite the historical
``intrinsic radius'' name; it is not the geometric radius or diameter of a
hyperbolic component.  At the period-two nucleus $c^*=-1$, $w_2=-2$ and
$d_2=-1$, giving $\mathrm{Scale}=1/2$ and $r_{\mathrm{int}}=2$.
The actual period-two bulb, whose cycle multiplier is $4(c+1)$, has
radius $1/4$ and diameter $1/2$, demonstrating
the distinction.  The reported scale is stored
as a \code{HighPrecision} value in \code{FeatureSummary} and can be used by
the renderer to set an appropriate zoom level or to determine the additional
bits of precision needed to render the component's interior.

\subsection{Checkpoint and resume}
\label{subsec:ff-checkpoint}

High-precision Newton--Raphson refinement can require minutes to hours
at extreme periods, because each outer step evaluates the critical orbit
for $p$~iterations at the current working precision.  To avoid losing
progress when the user aborts (\code{Ctrl+Alt}), closes the
application, or encounters an error, \FractalShark{} writes periodic
checkpoints to a file (\code{nr\_checkpoint.txt}) and resumes from the
most recent checkpoint on the next invocation.

\paragraph{Checkpoint contents.}
Each checkpoint records the current candidate~$c$ (real and imaginary
parts as decimal strings at full working precision), the detected
period~$p$, the MPIR precision in bits, the outer Newton iteration
index, and the inner-loop iteration count within the current $p$-step
evaluation, the orbit state $z$, first derivative, and reduced second
derivative.  The record also identifies the main, final-correction, or
complete phase, and retains the seed and search-radius information needed
for acceptance checks.  Storing the inner-loop position and its state allows resumption
\emph{mid-evaluation}, not only at outer-step boundaries.

\paragraph{Asynchronous writer.}
Converting high-precision \code{mpf\_t} values to decimal strings via
\code{mpf\_get\_str} and writing them to disk can take non-trivial
time at high precision.  A dedicated \code{CheckpointWriter} background
thread accepts \code{CheckpointSnapshot} objects (which own independent
copies of the \code{mpf\_t} data) and performs the serialization and
file I/O without blocking the Newton iteration.  The snapshot is
written to a temporary file and atomically renamed to the checkpoint
path, ensuring that the on-disk checkpoint is always in a consistent
state.

\paragraph{Save points.}
Checkpoints are written at two granularities:
\begin{itemize}
  \item After an outer Newton step completes and at final-correction phase
        transitions.
  \item Periodically during the inner-loop evaluation, via a progress
        callback that triggers whenever a configurable number of
        inner iterations have elapsed.
\end{itemize}
Both paths route through the same asynchronous writer, so at most one
background write is in flight at a time.

\paragraph{Resume.}
On entry to \code{RefinePeriodicPoint}, the code checks for an existing
checkpoint file.  If the file exists and its stored period and precision
match the current request, the saved candidate~$c$ and iteration state
are loaded and the refinement loop skips ahead to the checkpointed
position.  If the period or precision do not match (indicating a
different feature or a re-run at different settings), the checkpoint is
ignored and refinement starts from scratch.  The endpoint evaluators interpret
$\code{startIter}=0$ as fresh initialization, setting $z_0=d_0=q_0=0$.
A positive $\code{startIter}$ instead consumes all supplied checkpoint states
and performs only the remaining steps.  For example, resuming period $1000$
at index $600$ requires $400$ further steps at the same candidate $c$;
resetting the derivatives or orbit would evaluate a different state.  A
complete-phase checkpoint skips both main and final correction; a final-phase
checkpoint resumes only the final evaluation.

\paragraph{Abort safety.}
When the \code{AbortMonitor} sets the stop flag, the refinement loop
exits early \emph{after} writing a final checkpoint.  Progress is
therefore preserved: the user can later invoke \emph{Resume NR
Refinement} from the Feature Finder menu to continue from where the
computation was interrupted.

\subsection{Threading model}
\label{subsec:ff-threading}

\paragraph{Multithreaded CPU implementation.}
The dominant cost of the Feature Finder is the forward evaluation
\code{EvaluateCriticalOrbitAndDerivs}, which performs the inner loop
described in \cref{subsubsec:ff-inner} at high precision using GMP
\code{mpf\_t} arithmetic.  Because the 7~high-precision multiplies within
each iteration are independent (\cref{subsec:ff-inner-gpu}),
\FractalShark{} dispatches them to a pool of seven spin-based worker
threads (\code{NWORK~=~7})---one worker per multiply.
Synchronization uses cache-line-aligned (\code{alignas(64)}) generation
counters with atomic acquire/release semantics---no mutexes or condition
variables---so the main thread and workers coordinate with minimal latency.
The main thread composes the results after each batch completes and
advances to the next iteration.

A single-threaded fallback (\code{EvaluateCriticalOrbitAndDerivsST}) with
identical logic is also provided.

The Newton outer loop (\cref{subsubsec:ff-outer}) in
\code{FindPeriodicPoint\_Common} is inherently sequential: each Newton step
$c \leftarrow c - z_p / (\partial z_p / \partial c)$ depends on the
preceding evaluation's output, so successive steps cannot overlap.

GPU acceleration of the inner loop via the \code{HpSharkFloat} backend
is described in \cref{subsec:ff-inner-gpu}; this path is fully
implemented and dispatches through \code{Orbit\-Endpoint\-Evaluator} into the
fused Ref2/v2 kernel.

With the core iteration algorithms and precision infrastructure now
described, the next chapter turns to the post-iteration processing
pipeline---antialiasing, color mapping, and rendering orchestration---that
transforms raw escape-time counts into the final displayed image.
